\cotefile{37}
\scanpage{64}
% letterhead
\begin{center}
INSTITUT\\
DES HAUTES ÉTUDES SCIENTIFIQUES\\[-.3ex]
\rule{1.3cm}{.4pt}
\end{center}

{\reversemarginpar\margin{$\unsure{\alpha|\times}$ \unsure{App|AW}}}
\diagram{\(\begin{array}{r@{\quad}l}
\text{p. }43 & (\overline{\mathrm{I}\underline{\mathrm{X}}}\quad\ldots\quad)\\
\text{p. }47 & (\underline{\overline{\mathrm{IV}}}\qquad)\\
\text{p. }51 & (\underline{\overline{\mathrm{IV}}}\qquad)\\
\text{p. }52 & (\underline{\overline{\mathrm{IV}}}\qquad)\\
\text{p. }56 & (\underline{\overline{\mathrm{IX}}}\qquad)\\
\text{p. }59 & (\underline{\overline{\mathrm{IV}}}\qquad)\\
\text{p. }70 & (\mathrm{IX}\qquad)\\
\text{p. }\unsure{71|72|74} & (\underline{\overline{\mathrm{IX}}}\qquad)\\
\text{p. }75 & (\text{id}\qquad)
\end{array}\)}
% archive sheet number: 63

\scanpage{73}
\begin{flushright}
SGA 7
\end{flushright}
% archive sheet number: 72

\scanpage{78}
SGA 7 $\underline{\overline{\mathrm{VI}}}$ App

après 1.2 harmoniser notation $L\Omega_{\unsure{X/S|X/Y|X_{\varphi}|X/k}}$ pour compl. cot. rel.

mettre bibliographie \unsure{commune} à l'exposé et l'App.
% archive sheet number: 77

\scanpage{79}
% letterhead
\begin{center}
INSTITUT\\
DES HAUTES ÉTUDES SCIENTIFIQUES\\[-.3ex]
\rule{1.3cm}{.4pt}
\end{center}

\deleted{SGA 7 $\underline{\overline{\mathrm{VIII}}}$}

\deleted{Faux schémas abéliens (ou \unsure{Motifs|motifs} mixtes effectifs \unsure{de|des} poids 1)}

\deleted{Éléments de N.S. et correspondances symétriques.}

\noindent\rule{.62\linewidth}{.4pt}

$\underline{\overline{\mathrm{I}}}$ \quad Le groupe d'inertie $I$ a un $p$-ss-gpe de Sylow infini

Cf Serre \unsure{Corps claux|Corps locaux} \unsure{Coh. galois.|Cohom galois.}, p.144, lemme 4.

\noindent\rule{.88\linewidth}{.4pt}

$\underline{\overline{\unsure{\mathrm{VI}|\mathrm{VII}}}}$ \quad Traiter aussi le ``local invariant cycle problem'' pour \deleted{des \illegible{1}}\inserted{$H^{1}$}

\unsure{Signaler|sign.} la conjecture de Griffiths en dim. quelc.

\noindent\rule{.96\linewidth}{.4pt}
% archive sheet number: 78

\scanpage{80}
% letterhead
\begin{center}
INSTITUT\\
DES HAUTES ÉTUDES SCIENTIFIQUES\\[-.3ex]
\rule{1.3cm}{.4pt}
\end{center}

\diagram{\(\begin{array}{@{}l@{\;}r@{\quad}l@{\quad}l@{}}
\text{SGA 7 }\underline{\overline{\mathrm{IX}}} & \text{(iii)} & \text{l. }-5 & \text{\uline{Bibl}. ajouter [22 bis]}\\
 & 1 & \text{l. }5 & [13]\\
 & 33 & \text{l. }10 & [24]\\
 & 35 & \text{l. }5 &\\
 & 37 & \text{l. }8 &\\
 & 69 & \text{l. }-3 &\\
 & 70 & \text{l. }2 &\\
 & 73 & \text{l. }12 &\\
 & 87 & \text{l. }6 & \text{\unsure{livre} de Gabriel-Demazure}\\
 & 124 & \text{l. }-8 &\\
 & 128 & \multicolumn{2}{l}{\text{l. }4,\ \text{l. }-9}\\
 & 134 & \text{l. }1 &\\
 & 146 & \text{l. }5 &\\
 & 150 & \text{l. }-3 &\\
 & 157 & \text{l. }5 &\\
 & 1\deleted{\unsure{7|5}}59 & \multicolumn{2}{l}{\text{l. }11,\ \text{l. }-7}\\
 & 169 & \text{l. }4 & \text{deux fois}
\end{array}\)}
% archive sheet number: 79

\scanpage{81}
% letterhead
\begin{center}
INSTITUT\\
DES HAUTES ÉTUDES SCIENTIFIQUES\\[-.3ex]
\rule{1.3cm}{.4pt}
\end{center}

SGA 7 $\underline{\overline{\mathrm{XVII}}}$

\diagram{\(\begin{array}{ll}
\text{p. }2 & 1.\\
\text{p. }7 & \text{l. }-7\\
\text{p. }8 & \text{l. }17\\
\text{p. }24 & \text{l. }8
\end{array}\)}
% archive sheet number: 80

\scanpage{82}
% letterhead
\begin{center}
INSTITUT\\
DES HAUTES ÉTUDES SCIENTIFIQUES\\[-.3ex]
\rule{1.3cm}{.4pt}
\end{center}

Appendice à $\underline{\overline{\mathrm{XXI}}}$ de Katz : torsion en coh. cristalline

$\underline{\overline{\mathrm{XXII}}}$ \quad Référence à la démonstration de Deligne du th. de \unsure{congruences|congruence} de Katz (généralisée)

{\reversemarginpar\margin{$\underline{\overline{\mathrm{VI}}}$\\ou $\underline{\overline{\mathrm{XIII}}}$\unsure{?|'}}}Vérifier qu'on y trouve le fait qu'en déformant une singularité quadratique ordinaire, \unsure{on trouve|--- donne} au plus autant (y inclus cas de car. rés. 2 et dim relative paire). En d'autres termes, \unsure{l'ens|l'ens.} des pts où un morphisme plat \deleted{de}\inserted{loc.} \unsure{prés. finie} est lisse ou a une singularité quadratique ordinaire est \unsure{tjrs.|tjs.} \uline{ouvert}.
% archive sheet number: 81

\scanpage{83}
SGA 7 $\underline{\overline{\mathrm{XVIII}}}$ \qquad p. 10 l. 2 réf à \deleted{\unsure{SGA 7|SGA|Exp} $\underline{\overline{\unsure{\mathrm{XVIII}|\mathrm{XIII}}}}$ Exp} SGA 5 $\underline{\overline{\mathrm{IV}}}$

\diagram{\(\begin{array}{lll}
\text{p. }17 & \text{l. }9 & \text{réf \unsure{Exp|l'Exp}}\ \underline{\overline{\mathrm{XIV}}}\\
\text{p }25 & \text{l. }-2 & \text{---\quad SGA 5 }\underline{\overline{\mathrm{IV}}}\\
\text{p. }34 & \text{l }4 & \text{---}\quad\underline{\overline{\mathrm{XVI}}}\\
\text{p. }44 & \text{l\unsure{,|.} }1 &\\
\text{p }64 & \text{l. }-7,\ \text{l. }-2 & \text{--}\quad\underline{\overline{\mathrm{XVI}}}\\
\text{p. }66 & \text{l. }6 & \text{---}\quad\underline{\overline{\mathrm{XIII}}}\\
 & \text{l. }-5 & \text{id ?}\\
\text{p. }68 & \text{l. }4 & \underline{\overline{\mathrm{XVI}}}\\
 & \text{l. }-6 & \underline{\overline{\mathrm{XII}}}
\end{array}\)}
% archive sheet number: 82

\scanpage{84}
\begin{flushright}
\uline{Réédition EGA I}\\
\deleted{\unsure{et} \illegible{1}}
\end{flushright}
% archive sheet number: 83

\scanpage{95}
% letterhead
\begin{center}
INSTITUT\\
DES HAUTES ÉTUDES SCIENTIFIQUES\\[-.3ex]
\rule{1.3cm}{.4pt}
\end{center}

10.12. Stabilité par $\underset{\longleftarrow}{\lim}$ (\unsure{en un|en vue|en ex|ou non} ?)
% archive sheet number: 94

\scanpage{99}
\begin{flushright}
EGA $\underline{\overline{\mathrm{II}}}$
\end{flushright}
% archive sheet number: 98

\scanpage{100}
% letterhead
\begin{center}
INSTITUT\\
DES HAUTES ÉTUDES SCIENTIFIQUES\\[-.3ex]
\rule{1.3cm}{.4pt}
\end{center}

\textcircled{a} \quad \unsure{transférer} EGA $\underline{\overline{\mathrm{IV}}}$ 4.8
\[
\left\{\begin{aligned}
&\text{Également \unsure{suite|critère|écrit} des}\\
&\text{morphismes affines des}\\
&\text{\unsure{préschémas|préschèmes} formels.}
\end{aligned}\right.
\]

{\reversemarginpar\margin{!}}Appendice \deleted{\unsure{par|pour}} \unsure{Murre|Mémoire} aux schémas formels ?

\textcircled{b} \quad Cf EGA $\underline{\overline{\mathrm{II}}}$ 6.5.10 ajouter formules de la norme

6.5.11 ajouter formules pour le cas où $X'\unsure{\mathbin{-}|\to}X$ est gén. décomposé (du moins \unsure{géom!|géomt.})

6.5.\unsure{9|1|7} renvoi à EGA $\underline{\overline{\mathrm{IV}}}$ 21.\unsure{.}

\textcircled{c} \quad Remonter EGA $\underline{\overline{\mathrm{IV}}}$ 19.4.11

(revu et augmenté)

\textcircled{d} \quad 7.4.\unsure{2|3}. Admettre courbe vide ---
% archive sheet number: 99

